\title{FlashAttention-3: Fast and Accurate Attention with Asynchrony and Low-precision}

\begin{abstract}
As the foundational component within the widely adopted Transformer models, attention serves as a significant computational bottleneck, particularly in large language models and long-context tasks. \textsc{FlashAttention} introduced a method for accelerating attention on GPUs by reducing the volume of memory accesses. Nevertheless, its approach does not exploit advanced features available in state-of-the-art hardware, such as the H100 GPU, where \textsc{FlashAttention-2} manages to utilize only 35\% of the computational capacity. In this work, we propose three principal techniques to enhance attention performance on Hopper GPUs: (1) leveraging the asynchronous capabilities of Tensor Cores and TMA to concurrently schedule computation and memory transfer through warp specialization, (2) integrating block-level matrix multiplication and softmax stages via interleaving, and (3) employing block quantization with incoherent processing by capitalizing on native FP8 support in hardware. Our proposed framework, \textsc{FlashAttention-3}, achieves a performance boost of 1.5-2.0$\times$ on H100 GPUs—yielding FP16 throughput up to 740 TFLOPs/s (representing 75\% hardware utilization), and approaching 1.2 PFLOPs/s when utilizing FP8 precision. Furthermore, we show that FP8 \textsc{FlashAttention-3} attains numerical error that is 2.6$\times$ lower than that of standard FP8 attention implementations.
\end{abstract}

\section{Introduction}
\label{sec:intro}

Within the Transformer architecture~\citep{vaswani2017attention}, the primary source of computational expense arises from the attention mechanism, specifically due to the self-attention computation between queries and keys, which exhibits a quadratic growth with respect to the input sequence length. Extending the feasible attention span to accommodate longer sequences would enable advances such as more complex modeling and reasoning across several lengthy documents~\citep{guo2021longt5,shaham2022scrolls,peng2023yarn}, as well as handling multiple files in extensive code repositories~\citep{roziere2023code, li2023starcoder}. This would further facilitate the integration of additional modalities, including high-resolution imagery~\citep{chen2022scaling}, audio streams~\citep{gulati2020conformer}, and video data~\citep{ho2022video}, along with new use cases such as leveraging extensive user interaction history~\citep{sun2019bert4rec} or supporting agents working over long sequences of actions~\citep{yao2022react}. Consequently, substantial research efforts have been invested in accelerating the attention mechanism for long contexts, encompassing methods like approximate computation~\citep{katharopoulos2020transformers,choromanski2020rethinking, tay2020efficient}, optimizations at the software level~\citep{rabe2021self,dao2022flashattention,kwon2023efficient}, and the development of alternative neural architectures~\citep{peng2023rwkv,sun2023retentive,gu2023mamba}.

This study extends upon the contributions of \citet{dao2022flashattention}, who developed exact-attention methods informed by the GPU’s hardware architecture and execution behavior. Dao et al. \citep{dao2022flashattention} introduced the \textsc{FlashAttention} algorithm, which utilized an innovative tiling approach to execute attention calculations. By consolidating all attention-related operations within a singular GPU kernel, this design removed the necessity for intermediate accesses to slower global memory, thereby improving efficiency. In a subsequent work, \citet{dao2023flashattention2} presented \textsc{FlashAttention-2}, wherein the algorithm was redesigned to also parallelize along the sequence length axis and to process key and value matrices blockwise within the inner loop of the forward pass, resulting in more effective workload balancing and resource allocation on GPUs. Despite these advances, our analysis indicates that \textsc{FlashAttention-2} attains suboptimal hardware utilization compared to highly optimized matrix-multiplication (GEMM) kernels, with utilization rates around 35\% in contrast to the 80-90\% observed on the Hopper H100 GPU. One plausible explanation is implementation-specific factors, such as a failure to leverage Hopper’s dedicated instructions rather than defaulting to Ampere instructions when using Tensor Cores. Previous investigations, including ThunkerKitten~\citep{spector2024thunder} and cuDNN 9~\citep{cudnn9}, have demonstrated that incorporating Hopper-oriented instructions alongside tile-centric programming abstractions can both accelerate attention computations and streamline code implementation.

At a deeper level, the algorithm underlying \textsc{FlashAttention-2} operates under a straightforward synchronous paradigm, explicitly omitting the incorporation of asynchrony and reduced-precision strategies in its architecture. Asynchrony emerges from the deployment of hardware accelerators purpose-built for critical machine learning tasks: for instance, dedicated processing units such as Tensor Cores (for matrix multiplication) or Tensor Memory Accelerator (TMA) for managing data movement, both functioning independently of general CUDA cores tasked with other computational roles—ranging from logic and integer to floating-point operations. The adoption of low-precision representations—FP8 in Hopper and FP4 in Blackwell, following predecessors like FP16 (introduced with Pascal in 2017) and BF16 (debuted in Ampere in 2020)—continues to enable dramatic gains in throughput, achieving twofold or even fourfold speedups without increasing energy consumption or chip area. We discuss the enhancements provided by Hopper in these respects in \cref{subsec:hardware}.  
The primary technical hurdle is the reengineering of \textsc{FlashAttention-2} to harness these advanced hardware functionalities: employing asynchrony necessitates effective overlap of the computation phases for matmul and softmax, despite inherent data dependencies; leveraging low-precision formats requires meticulous design to suppress quantization errors, which is particularly critical for rare or extreme activations present in LLMs~\citep{dettmers2208llm, sun2024massive}.

With this objective, we introduce \textsc{FlashAttention-3}, which integrates and advances three innovative concepts to enhance efficiency on the latest GPU platforms:\footnote{Our findings are contextualized within the framework of NVIDIA’s Hopper architecture; nevertheless, the proposed algorithm remains applicable to any GPU platform equipped with strong asynchronous processing and support for low-precision operations.}

\begin{enumerate}

\item \textbf{Producer-Consumer asynchrony:} We introduce a warp-specialized software pipelining strategy that leverages asynchronous execution between data movement and Tensor Cores. By separating the roles of producers and consumers into distinct warps, the algorithm can more effectively mask both memory access and instruction issue latencies.
\item \textbf{Hiding softmax under asynchronous block-wise GEMMs:} We mask the latency of non-GEMM softmax-related operations (such as floating point multiply-add and exponentiation), which typically offer lower throughput, by overlapping them with asynchronous WGMMA GEMM instructions. To achieve this, we modify the \textsc{FlashAttention-2} workflow to remove specific sequential constraints between softmax computation and GEMM operations. For instance, within the two-stage algorithm, while the softmax function processes a block from the scores matrix, the asynchronous WGMMA pathway can simultaneously prepare and compute the following block.
\item \textbf{Hardware-accelerated low-precision GEMM:} We redesign the forward computation to enable usage of FP8 Tensor Cores during GEMM, nearly achieving a twofold increase in effective TFLOPs/s. This adaptation entails handling the divergent memory layout requirements imposed by WGMMA, specifically regarding the in-memory arrangement of FP32 accumulator matrices versus FP8 operand blocks. To address accuracy degradation arising from reduced precision, we incorporate techniques such as block quantization and incoherent processing.

To empirically assess our approach, we perform benchmarks of \textsc{FlashAttention-3} on the H100 SXM5 GPU across a variety of parameter settings. Our results demonstrate that (1) FP16 provides a 1.5-2.0$\times$ acceleration in the forward computation over \textsc{FlashAttention-2} (achieving up to 740 TFLOPs/s), and achieves a 1.5-1.75$\times$ gain in the backward computation, (2) FP8 approaches a throughput of 1.2 PFLOPs/s, and (3) at large sequence lengths, FP16 achieves superior performance, while FP8 remains competitive\footnote{Specifically, when the head dimension is 64, \textsc{FlashAttention-3} FP8 exceeds the baseline, whereas for head dimensions of 128 and 256, its performance matches the baseline in non-causal scenarios and trails it in settings with causal masking.} compared to the latest attention implementation provided by NVIDIA’s cuDNN library. Additionally, we confirm that FP16 \textsc{FlashAttention-3} maintains numerical error on par with \textsc{FlashAttention-2} and yields improved precision compared to conventional attention kernels, due to retention of intermediate values (such as softmax scaling) in FP32 format. Furthermore, we observe that FP8 \textsc{FlashAttention-3}, when using block quantization and incoherent processing, attains 2.6$\times$ higher accuracy relative to standard attention employing per-tensor quantization, particularly in instances exhibiting outlier features.

We have released \textsc{FlashAttention-3} under a permissive open-source license\footnote{\textsc{FlashAttention-3}
  is available at \texttt{https://github.com/Dao-AILab/flash-attention}}. Additionally, we intend to incorporate it into the PyTorch and Hugging Face ecosystems to ensure broad accessibility for both researchers and practitioners.

\section{Background: Multi-Head Attention and GPU Characteristics}
\label{sec:background}

\subsection{Multi-Head Attention}
\label{subsec:multi_head_attn}

Let $\mathbf{Q}, \mathbf{K}, \mathbf{V} \in \mathbb{R}^{N \times d}$ be the query, key and value input sequences associated to a single head, where $N$ is the sequence length and $d$ is the head dimension. Then the attention output $\mathbf{O}$ is computed as:

\begin{equation*}
  \mathbf{S} = \alpha \mathbf{Q} \mathbf{K}^\top \in \mathbb{R}^{N \times N}, \quad \mathbf{P} = \mathrm{softmax}(\mathbf{S}) \in \mathbb{R}^{N \times N}, \quad \mathbf{O} = \mathbf{P}\mathbf{V} \in \mathbb{R}^{N \times d},
\end{equation*}

In this formulation, the $\mathrm{softmax}$ operation is computed across each row, with the scaling factor commonly chosen as $\alpha = 1/\sqrt{d}$. To address potential numerical instability caused by the exponential function, it is standard practice to subtract $\mathrm{rowmax}(\mathbf{S})$ from $\mathbf{S}$. In the context of multi-head attention (MHA), distinct query, key, and value projections are employed for each head. This allows the calculations to be executed in parallel over both multiple heads and batches, resulting in the complete output tensor.

Now let $\phi$ be a scalar loss function and let $\mathbf{d}(-) = \partial \phi / \partial (-)$ be notation for the gradient. Given the output gradient $\mathbf{dO} \in \mathbb{R}^{N \times d}$, we compute $\mathbf{dQ}$, $\mathbf{dK}$, and $\mathbf{dV}$ according to the chain rule as follows:

\begin{align*}
  \mathbf{dV} &= \mathbf{P}^\top \mathbf{dO} \in \mathbb{R}^{N \times d} \\
  \mathbf{dP} &= \mathbf{dO} \mathbf{V}^\top \in \mathbb{R}^{N \times N} \\
  \mathbf{dS} &= \mathrm{dsoftmax} (\mathbf{dP}) \in \mathbb{R}^{N \times N} \\
  \mathbf{dQ} &= \alpha \mathbf{dS} \mathbf{K} \in \mathbb{R}^{N \times d} \\
  \mathbf{dK} &= \alpha \mathbf{dS}^\top \mathbf{Q} \in \mathbb{R}^{N \times d},
\end{align*}

In this context, the relationship $\mathbf{d}s = (\mathrm{diag}(p) - p p^\top)\mathbf{d}p$ holds, where $p = \mathrm{softmax}(s)$ defines $p$ as a vector-valued function of $s$. We denote the application of this derivative row-by-row as $\mathrm{dsoftmax}(\mathbf{dP})$. Moreover, this process can be efficiently parallelized over all heads and batch elements during the backward propagation phase of MHA.

\subsection{GPU hardware characteristics and execution model}
\label{subsec:hardware}

We discuss the GPU execution model features pertinent to \textsc{FlashAttention-3}, concentrating on the NVIDIA Hopper architecture as a specific realization of this model.

\paragraph{Memory hierarchy:} The GPU features a tiered arrangement of memory types, where the highest levels offer lower bandwidth and larger capacity (\cref{tab:gpu-hierarchy})\footnote{\citet{luo2024benchmarking} notes that shared memory achieves a bandwidth of 128 bytes per clock cycle for each SM; with 132 SMs operating at a boost frequency of 1830 MHz, this value is appropriately scaled.}. Global memory (GMEM), or HBM, resides off-chip and is shared across all streaming multiprocessors (SMs). Data stored in GMEM is automatically buffered in a dedicated L2 on-chip cache. Further down the hierarchy, each SM possesses its own segment of highly banked, user-controllable shared memory (SMEM) on-chip. At the lowest level, each SM maintains a private register file.

\paragraph{Thread hierarchy:} In the GPU programming paradigm, threads are structured into several logical layers. These range from individual threads to larger groupings: warps, which consist of 32 threads; warpgroups, each made up of 4 adjacent warps; threadblocks (also known as cooperative thread arrays or CTAs); threadblock clusters (a feature introduced in Hopper architecture); and, at the highest level, grids.

The interdependence between these two hierarchies is significant. Threads belonging to a single CTA are scheduled together on a single SM, while CTAs residing within a common cluster are concurrently scheduled on the same GPC. All threads within a CTA have direct access to SMEM, in contrast to RMEM, where each thread possesses up to 256 private registers accessible only to that thread.

\paragraph{Asynchrony and warp-specialization:}

GPUs achieve high throughput by leveraging parallelism and asynchronous operations to mask delays associated with memory access and computation. In Hopper architecture, the Tensor Memory Accelerator (TMA) serves as a specialized hardware component for facilitating asynchronous data transfers between GMEM and SMEM \cite[\S7.29]{cuda}. Moreover, in contrast to previous generations like Ampere, Hopper's Tensor Core—accessible through the warpgroup-wide WGMMA instruction \cite[\S9.7.14]{ptx}—operates asynchronously and is capable of reading inputs directly from shared memory.

By incorporating hardware mechanisms for asynchrony, warp-specialized kernels become feasible; within a CTA, warps can be distinctly assigned to either producer or consumer functions, performing exclusively data movement or computation tasks, respectively. This separation generally enhances the compiler’s capacity to produce more efficient instruction schedules \citep{warp-specialization-2011}. Furthermore, Hopper provides the capability for registers to be reallocated dynamically among warpgroups by using \verb|setmaxnreg| \cite[\S9.7.17.1]{ptx}. Consequently, warps engaged in MMAs can be provisioned with a greater portion of RMEM than those restricted to TMA operations, for which just a single thread is sufficient.

\paragraph{Low-precision number formats:}
\label{sec:low-precision-gpu}
State-of-the-art GPUs incorporate dedicated hardware specifically designed to speed up computations using low-precision formats. As an illustration, the WGMMA instruction is capable of utilizing the FP8 Tensor Cores within Hopper architecture, achieving twice the throughput per SM relative to using FP16 or BF16 formats.

To utilize FP8 WGMMA accurately, it is essential to recognize the specific data layout requirements for its input matrices. Consider the scenario of performing a GEMM operation to compute $A \times B^{\top}$, where $A$ is of shape $M\times K$ and $B$ is $N\times K$. Here, an operand is described as \emph{mn-major} when it is stored contiguously along the outer $M$ or $N$ dimensions, and as \emph{k-major} when contiguous storage is along the inner $K$ dimension. For FP16 WGMMA, SMEM operands can be provided in either mn-major or k-major configurations. In contrast, FP8 WGMMA only permits the k-major layout for input operands. This becomes particularly problematic in applications like attention mechanisms, where combining consecutive GEMM operations in a unified kernel is desired; the difference between the layouts of FP32 accumulators and FP8 operands hinders the chaining of dependent FP8 WGMMA calls.

Regarding attention mechanisms, these layout constraints necessitate specific adjustments to the construction of an FP8 algorithm, as detailed in \cref{sec:algofp8}.

\subsection{Standard Attention and Flash Attention} In line with the framework established by \citet{dao2022flashattention}, we use the term \textbf{standard attention} to refer to a GPU-based attention implementation that writes the intermediary matrices $\mathbf{S}$ and $\mathbf{P}$ to high-bandwidth memory (HBM). \textsc{FlashAttention} introduces a critical improvement: by applying a local softmax reduction, it eliminates costly intermediate memory accesses and instead combines the entire attention computation within a single kernel invocation. Within the main consumer loop described in \cref{alg:flash3_wgmma_ws_only}, the local softmax is carried out during lines \ref{code-ws:softmax_start}-\ref{code-ws:softmax_end}, along with the rescaling of $\mathbf{O}$ in blocks. A straightforward argument demonstrating that this method yields the correct computation for $\mathbf{O}$ is provided in \cite[\S 2.3.1]{dao2023flashattention2}.

\section{FlashAttention-3: Algorithm}
\label{sec:algo}

In this section, we present an overview of the \textsc{FlashAttention-3} algorithm. Our primary focus is on the forward pass, while details concerning the backward pass can be found in~\cref{sec:algo_ws_bwd}. We begin by outlining the process of incorporating warp-specialization and a circular SMEM buffer within the foundational framework established by \textsc{FlashAttention-2}. Subsequently, we discuss leveraging WGMMA asynchrony to establish a two-stage pipeline in which the GEMM and softmax computations can be overlapped. Lastly, we outline the adaptations necessary to support FP8, specifically with respect to layout compatibility and precision, utilizing block quantization and methods for managing incoherent processing.

\subsection{Producer-Consumer asynchrony through warp-specialization and pingpong scheduling}
\label{sec:algo_ws}

\paragraph{Warp-specialization}
Much like in \textsc{FlashAttention-2}, the forward computation in \textsc{FlashAttention-3} can be executed with high parallelism across the batch dimension, number of attention heads, and the length of the query sequence. Consequently, a CTA-level abstraction suffices for exposition: this perspective processes a tile $\mathbf{Q}_i$ extracted from the query matrix and produces the associated tile $\mathbf{O}_i$ in the output. For clarity, the initial description focuses on a warp-specialized strategy that utilizes a circular SMEM buffer, but does not incorporate the concurrent execution of GEMM and softmax computations. Given head dimension $d$ and sequence length $N$, we select a query block size $B_r$, segmenting $\mathbf{Q}$ into $T_r = \lceil \frac{N}{B_r} \rceil$ partitions: $\mathbf{Q}_1, .., \mathbf{Q}_{T_r}$.

\begin{algorithm}[H]
    \caption{\small\label{alg:flash3_wgmma_ws_only}\textsc{FlashAttention-3} forward pass \textbf{excluding} intra-consumer overlap -- CTA perspective}
    \begin{algorithmic}[1]
\REQUIRE Input matrices $\mathbf{Q}_i \in \mathbb{R}^{B_r \times d}$, $\mathbf{K}$, and $\mathbf{V} \in \mathbb{R}^{N \times d}$ stored in HBM; key block size $B_c$ with the number of chunks $T_c = \lceil \frac{N}{B_c} \rceil$.
\STATE Set up pipeline structure to coordinate barriers using a $s$-stage circular SMEM buffering scheme.
\IF {executing within the producer warpgroup}
\STATE Release a predefined count of registers.
\STATE Load $\mathbf{Q}_i$ from HBM into shared memory.
\STATE Once load completes, signal to consumers that $\mathbf{Q}_i$ is ready.
\FOR{$0 \le j < T_c$}
    \STATE Wait for the consumer to finish with the $(j\,\%\,s)$th buffer stage.
    \STATE Trigger HBM-to-shared memory transfers for $\mathbf{K}_j, \mathbf{V}_j$ into the $(j\,\%\,s)$th buffer slot.
    \STATE After finishing the transfers, signal readiness of $\mathbf{K}_j, \mathbf{V}_j$ to consumers.
\ENDFOR
\ELSE \STATE Claim a preassigned number of registers, scaling with the number of consumer warps.
\STATE Initialize on-chip buffers: $\mathbf{O}_i = (0) \in \mathbb{R}^{B_r \times d}$ and scalar vectors $\ell_i, m_i = (0), (-\infty) \in \mathbb{R}^{B_r}$.
\STATE Wait for notification that $\mathbf{Q}_i$ is available in shared memory.
\FOR{$0 \le j < T_c$}
\STATE Wait until $\mathbf{K}_j$ is accessible in shared memory.
\STATE Calculate $\mathbf{S}_i^{(j)} = \mathbf{Q}_i \mathbf{K}_j^T$ using SS-GEMM. Commit operation and synchronize. \label{alg:ws_only_gemm1}
\STATE Record $m_i^{\mathrm{old}} = m_i$; update $m_i = \mathrm{max}(m_i^{\mathrm{old}}, \mathrm{rowmax}(\mathbf{S}_i^{(j)}))$. \label{code-ws:softmax_start}
\STATE Evaluate $\widetilde{\mathbf{P}}_i^{(j)} = \mathrm{exp}(\mathbf{S}_i^{(j)} - m_i)$ and update $\ell_i = \mathrm{exp}(m_i^{\mathrm{old}} - m_i) \ell_i + \mathrm{rowsum}(\widetilde{\mathbf{P}}_i^{(j)})$. \label{code-ws:softmax_end}
\STATE Pause execution until $\mathbf{V}_j$ is ready in shared memory.
\STATE Update $\mathbf{O}_i = \mathrm{diag}(\mathrm{exp}(m_i^{\mathrm{old}} - m_i))^{-1} \mathbf{O}_i + \widetilde{\mathbf{P}}_i^{(j)} \mathbf{V}_j$ using RS-GEMM. Commit and synchronize. \label{alg:ws_only_gemm2}
\STATE Mark the $(j\,\%\,s)$th buffer slot as released to the producer.
\ENDFOR
\STATE Finalize $\mathbf{O}_i = \mathrm{diag}(\ell_i)^{-1} \mathbf{O}_i$ and compute $L_i = m_i + \log(\ell_i)$.
\STATE Store $\mathbf{O}_i$ and $L_i$ into HBM as the $i$th output block in $\mathbf{O}$ and $L$.
\ENDIF
\end{algorithmic}
\end{algorithm}

In our realization of \cref{alg:flash3_wgmma_ws_only} on Hopper, memory (de)allocations are managed via \verb|setmaxnreg|, while data for $\mathbf{Q}_i$ and $\{ \mathbf{K}_j, \mathbf{V}_j \}_{0 \leq j < T_c}$ are brought in using TMA loads. The consumer mainloop uses WGMMA instructions to perform GEMMs, where the SS or RS prefix designates whether shared memory or registers provide the first operand. To understand the control flow in \cref{alg:flash3_wgmma_ws_only}, it is important to recognize that asynchronous TMA loads can be issued independently and do not block one another. Additionally, the producer mainloop initiates its first $s$ iterations without issuing any waits, as these initial passes are dedicated to populating the buffer.

\paragraph{Pingpong scheduling} The inherent asynchrony between WGMMA and TMA, when combined with warp specialization, enables the possibility of interleaving the softmax calculation of one warpgroup with the GEMM performed by another. This overlapping can be beneficial, given the significantly higher throughput of matmul instructions over non-matmul operations on contemporary hardware accelerators. For instance, on the H100 SXM5 GPU, FP16 matrix multiplications can reach 989 TFLOPS, whereas special functions such as the exponential (essential for softmax) achieve only 3.9 TFLOPS\footnote{The CUDA programming guide notes that each streaming multiprocessor (SM) can execute 16 special function operations per cycle. Multiplying 16 by the total number of SMs (132) and the clock frequency (1830 MHz) yields 3.9 TFLOPS for special functions.}. Considering an FP16 attention forward pass with a head size of 128, the number of FLOPS for matmul is 512 times that of exponentials, but the hardware supports exponentials at only 1/256 of the throughput. As a result, exponential calculations may occupy as much as 50\% of the compute cycles relative to matmuls. This disparity becomes even more pronounced in FP8, where matmul speed doubles while exponential throughput does not improve.

Given that the computation of the exponential function is handled by a distinct hardware component (the multi-function unit), optimal efficiency is achieved when this operation overlaps with the matmul being processed on the Tensor Cores. To orchestrate this concurrency, we implement synchronization barriers (\texttt{bar.sync} instructions) that ensure the GEMM operations (specifically, GEMM1 -- $\mathbf{P} \mathbf{V}$ for one iteration, and GEMM0 -- $\mathbf{Q} \mathbf{K}^\top$ for the subsequent iteration) for warpgroup 1 are completed ahead of those for warpgroup 2. This mechanism schedules the softmax computation for warpgroup 1 to coincide with the execution of GEMMs by warpgroup 2. Subsequently, the workflow alternates—warpgroup 2 proceeds with softmax calculations while warpgroup 1 handles the next set of GEMMs. This interleaved or ``pingpong’’ scheduling approach is visually represented in~\cref{fig:pingpong_scheduling}. While in practice this pingpong scheduling does not always precisely match the depicted execution order, empirical results demonstrate a notable performance gain (for instance, increasing from 570 TFLOPS to a range of 620--640 TFLOPS in FP16 forward runs with a head dimension of 128 and sequence length of 8192).

\paragraph{Attention variants} For both multi-query attention~\citep{shazeer2019fast} and grouped query attention~\citep{ainslie2023gqa}, our method adheres to the strategy implemented in \textsc{FlashAttention-2}, modifying the tensor indexing so as to prevent redundant copies of $\mathbf{K}$ and $\mathbf{V}$ within HBM.

\subsection{Intra-warpgroup overlapping GEMMs and softmax}

It is possible to interleave certain operations from the softmax with some operations from the GEMMs, even inside a single warpgroup. Below, we outline an approach to achieve this overlap.

In the attention algorithm, operations within the inner loop (main loop) have sequential dependencies that impede parallelization within a single iteration. For example, (local) softmax (lines \ref{code-ws:softmax_start} to \ref{code-ws:softmax_end}) relies on the output $\mathbf{S}_i^{(j)}$ of the first GEMM, while the second GEMM takes its result $\widetilde{\mathbf{P}}_i^{(j)}$ as an operand. Indeed, the wait statements in lines \ref{alg:ws_only_gemm1} and \ref{alg:ws_only_gemm2} of \cref{alg:flash3_wgmma_ws_only} serialize the execution of softmax and GEMMs. However, we can break these dependencies by pipelining across iterations through additional buffers in registers. Pursuing this idea, we propose the following two-stage\footnote{Note that the number of stages of the overlapping scheme is bounded by, but need not equal, the number $s$ of stages in the circular SMEM buffer.} GEMM-softmax pipelining algorithm:

\begin{algorithm}[H]
    \caption{\small\label{alg:flash3_wgmma}\textsc{FlashAttention-3} Consumer Warpgroup Forward Computation}
    \begin{algorithmic}[1]
      \REQUIRE  Input matrices $\mathbf{Q}_i \in \mathbb{R}^{B_r \times d}$, $\mathbf{K}, \mathbf{V} \in \mathbb{R}^{N \times d}$ residing in HBM, with key block dimension $B_c$ and number of key blocks $T_c = \lceil \frac{N}{B_c} \rceil$.
      \STATE Dynamically assign a predetermined number of registers, conditioned on the count of consumer warps.
      \STATE On-chip memory is used to initialize $\mathbf{O}_i = (0) \in \mathbb{R}^{B_r \times d}$, while $\ell_i, m_i$ are initialized as $(0), (-\infty) \in \mathbb{R}^{B_r}$, respectively.
      \STATE Synchronize until $\mathbf{Q}_i$ and the initial key block $\mathbf{K}_0$ are available in shared memory.
      \STATE Using WGMMA, perform $\mathbf{S}_{\mathrm{cur}} = \mathbf{Q}_i \mathbf{K}_0^T$. After committing, await completion.
      \STATE Free the initial buffer stage for $\mathbf{K}$.
      \STATE Compute $m_{i}$, $\tilde{\mathbf{P}}_{\mathrm{cur}}$, and $\ell_{i}$ from $\mathbf{S}_{\mathrm{cur}}$, and apply the necessary scaling to $\mathbf{O}_i$.
      \FOR{$1 \le j < T_c - 1$}
        \STATE Stall until block $\mathbf{K}_j$ is loaded to shared memory.
        \label{code:mainloop_start}
        \STATE Utilize WGMMA to calculate $\mathbf{S}_{\mathrm{next}} = \mathbf{Q}_i \mathbf{K}_{j}^T$. Commit this operation but proceed without waiting. \label{code:first_wgmma}
        \STATE Wait until the previous value block $\mathbf{V}_{j-1}$ is accessible in shared memory.
        \STATE Update $\mathbf{O}_{i} = \mathbf{O}_{i} + \tilde{\mathbf{P}}_{\mathrm{cur}} \mathbf{V}_{j-1}$ using WGMMA; submit for processing but do not pause for completion. \label{code:second_wgmma}
        \STATE Pause execution until completion of $\mathbf{Q}_i \mathbf{K}_{j}^T$ via WGMMA.
        \STATE Update $m_{i}$, $\tilde{\mathbf{P}}_{\mathrm{next}}$, and $\ell_{i}$ based on the computed $\mathbf{S}_{\mathrm{next}}$. \label{code:softmax}
        \STATE After $\tilde{\mathbf{P}}_{\mathrm{cur}} \mathbf{V}_{j-1}$ WGMMA completion, re-apply scaling to $\mathbf{O}_i$.
        \STATE Free the buffer slot $(j\,\%\,s)$ for $\mathbf{K}$ and $(j-1\,\%\,s)$ for $\mathbf{V}$, respectively.
        \STATE Move $\mathbf{S}_{\mathrm{next}}$ into $\mathbf{S}_{\mathrm{cur}}$ for subsequent iterations.
        \label{code:mainloop_end}
      \ENDFOR \STATE Hold until $\mathbf{V}_{T_c - 1}$ becomes available in shared memory.
      \STATE Finally, compute
      $\mathbf{O}_{i} = \mathbf{O}_{i} + \tilde{\mathbf{P}}_{\mathrm{last}} \mathbf{V}_{T_c - 1}$ via WGMMA, ensuring to commit the operation and wait for its conclusion.

\STATE Epilogue: Adjust $\mathbf{O}_{i}$ according to $m_{i}$. Derive $L_{i}$ utilizing both $m_{i}$ and $\ell_{i}$. Store $\mathbf{O}_{i}$ and $L_{i}$ into HBM as the $i$-th segment of $\mathbf{O}$ and $L$, respectively.
\end{algorithmic}
\end{algorithm}

\cref{alg:flash3_wgmma} serves as a substitute for the consumer component in \cref{alg:flash3_wgmma_ws_only}, thus forming the full \textsc{FlashAttention-3} algorithm in the context of FP16 computation. Conceptually, we treat WGMMA as synonymous with an asynchronous GEMM operation. In the main computational loop (see lines \ref{code:mainloop_start} to \ref{code:mainloop_end}), the execution of the second WGMMA instance in iteration $j$ (refer to line \ref{code:second_wgmma}) is scheduled concurrently with the softmax computations from the subsequent iteration $j+1$ (line \ref{code:softmax}), thereby enabling overlap of these operations.

Although the pipelined approach depicted earlier can yield significant performance improvements in theory, various practical challenges arise:

\paragraph{Compiler reordering} Although the pseudocode outlines an optimal sequence of execution, the compiler (NVCC) frequently reorders instructions during optimization. Such reordering may interfere with the intended arrangement of WGMMA and non-WGMMA instructions, thereby impacting the efficiency of the pipeline and potentially causing suboptimal performance or incorrect behavior. Examination of the generated SASS code confirms that the compiler introduces the expected overlap (refer to Section~\ref{sec:2-stage-sass}).

\paragraph{Register pressure} For peak efficiency, it is crucial to limit register spilling. Nevertheless, implementing the 2-stage pipeline increases the number of required registers, as it must preserve intermediate results and facilitate context switching between stages. For instance, each threadblock must retain an additional $\mathbf{S}_{\mathrm{next}}$ in registers, imposing an extra overhead equivalent to $B_r \times B_c \times \text{sizeof}(\text{float})$. The increased consumption of registers can limit the feasibility of larger block sizes—a commonly used optimization—due to their additional demand for registers. In practice, one must carefully weigh these factors, guided by profiling data.

\paragraph{3-stage pipelining} Building on the earlier 2-stage formulation, we introduce a 3-stage pipeline, which seeks to concurrently schedule the second WGMMA alongside the softmax phase. This structure theoretically improves Tensor Core usage even further, but at the cost of heightened register usage, since the pipeline's deeper structure requires the retention of more intermediate state. Consequently, the balance between maximizing tile size and pipeline depth becomes increasingly challenging. For an in-depth exposition and quantitative results concerning this 3-stage pipeline, see ~\cref{sec:3-stage}.

\subsection{Low-precision with FP8}
\label{sec:algofp8}

\textbf{Efficiency: layout transformations.} Executing the forward pass of \textsc{FlashAttention-3} using FP8 precision introduces extra difficulties related to layout compatibility that are absent when working with FP16.

To begin with, it is important to highlight that the input tensors $\mathbf{Q}$, $\mathbf{K}$, and $\mathbf{V}$ are usually stored such that the head dimension is contiguous in memory. However, meeting the k-major requirement of FP8 WGMMA for the second GEMM operation necessitates that $\mathbf{V}$—more specifically, the tiles of $\mathbf{V}$ stored in shared memory (SMEM)—be contiguous in the sequence length dimension. As the TMA load does not allow for changing which dimension is contiguous, we are left with two approaches: (1) pre-transpose $\mathbf{V}$ in global memory (GMEM) before computation, or (2) perform an in-kernel tile-wise transpose of $\mathbf{V}$ after tiles have been loaded into SMEM. For the first strategy, one can choose between (1a) integrating the transpose operation with the epilogue of a previous stage, such as the rotary embedding step, or (1b) launching a dedicated transpose kernel beforehand\footnote{A highly optimized transpose kernel can reach device-level bandwidth efficiency~\citep{colfax_cutlass_transpose_2024}.}, thereby swapping the strides between sequence length and head dimensions. Nonetheless, (1a) presents challenges for inclusion in standard library designs, and (1b) proves inefficient due to additional memory traffic, especially during inference, where memory bandwidth is a significant limitation.

Rather than alternative approaches, we utilize option (2) for FP8 \textsc{FlashAttention-3}. Within the in-kernel transpose phase, we exploit the capabilities of the LDSM (\verb|ldmatrix|) and STSM (\verb|stmatrix|) instructions. These allow a warp of threads to collaboratively transfer data between SMEM and RMEM in chunks of 128 bytes.\footnote{Although LDSM/STSM are specified in the PTX manual as operations for copying $8 \times 8$ matrices with 16-bit elements \cite[\S 9.7.13.4.15-16]{ptx}, pairing 8-bit elements enables the use of LDSM/STSM under FP8 settings. However, the transpose variants of these instructions cannot individually separate packed 8-bit values, making intermediate register shuffling necessary between LDSM and STSM to achieve the required tile-level transpositions; we leave out the specific implementation details here.} LDSM and STSM are highly efficient with registers, facilitating their execution within the producer warpgroup and permitting layout transpositions during memory copy operations. Additionally, after the initial iteration, it becomes possible to pipeline the transpose operation of the upcoming $\mathbf{V}$ tile during the concurrent execution of the two WGMMAs associated with the previous $\mathbf{V}$ tile and the current $\mathbf{K}$ tile.

Second, we note a distinction between FP8 WGMMA operations and those using FP16, specifically regarding the memory layout of the FP32 accumulator compared to that expected for operand A when stored in registers. Illustrative examples of these two layouts can be found in \cref{fig:rmem_accum} and \cref{fig:rmem_operand}, where the register assignments per thread are presented in the given order. By utilizing byte permutation instructions, it becomes possible to reformat the output of the first WGMMA’s accumulator into an arrangement that is compatible for input to a subsequent WGMMA, as well as with the layout of the $\mathbf{V}$ tile after the in-kernel transpose. Concretely, as shown in \cref{fig:rmem_accum}, the ordering is adjusted to
$$\{ \verb|d0 d1 d4 d5 d2 d3 d6 d7| \},$$
and this reordering of register values is repeated for every 8 bytes segment. This modification results, for the logical layout of the $\mathbf{P}$ tile, in a reorganization of its columns (for instance, columns $0189$ are mapped as the initial four columns). In order to ensure that the WGMMA operation produces the intended output tile, the in-kernel transpose can be designed to emit the $\mathbf{V}$ tile with a corresponding row reordering.\footnote{The flexibility provided by in-kernel transpose removes the need for shuffle instructions to transfer register values across threads, an approach previously addressed in~\citep{colfax_fp8_flashattention_2024}.}

\textbf{Accuracy: block quantization and incoherent processing.} The FP8 (e4m3) format allocates only 3 bits to represent the mantissa and 4 bits for the exponent, which results in a higher degree of numerical error relative to FP16/BF16 formats. Additionally, large-scale models are known to have outlier values~\citep{dettmers2208llm, sun2024massive} that are considerably larger than the bulk of values, thereby complicating the quantization procedure. Common practice is to apply per-tensor scaling~\citep{micikevicius2022fp8}, assigning a single scaling factor to each tensor (such as one for $\mathbf{Q}$, one for $\mathbf{K}$, and one for $\mathbf{V}$). To address the amplified numerical errors found in FP8 attention, we utilize two strategies:

\begin{enumerate}

\item \textbf{Block quantization}: A scalar is retained for each block, meaning that for $\mathbf{Q}$, $\mathbf{K}$, and $\mathbf{V}$, the tensor is partitioned into blocks of dimension $B_r \times d$ or $B_c \times d$, with quantization applied independently to each block. This quantization procedure can be combined with preceding operations, such as rotary embedding, without incurring extra runtime overhead, given that rotary embedding is limited by memory bandwidth. Since the \textsc{FlashAttention-3} algorithm is inherently structured around block-wise operations, it is feasible to apply a scaling factor to each block of $\mathbf{S}$ to compensate for block quantization, with negligible computational overhead.

\item \textbf{Incoherent processing}: To mitigate the impact of outlier values, $\mathbf{Q}$ and $\mathbf{K}$ are premultiplied by a random orthogonal matrix $\mathbf{M}$ prior to FP8 quantization. Since $\mathbf{M}$ satisfies $\mathbf{M} \mathbf{M}^\top = I$, it follows that $(\mathbf{Q} \mathbf{M}) (\mathbf{K} \mathbf{M})^\top = \mathbf{Q} \mathbf{K}^\top$, ensuring that the use of $\mathbf{M}$ does not affect the resultant attention. This process helps to redistribute outliers since each entry in $\mathbf{Q} \mathbf{M}$ or $\mathbf{K} \mathbf{M}$ becomes a random weighted sum of original entries, thereby decreasing the quantization error. Following the approach in \citet{chee2024quip} and~\citet{tseng2024quip}, $\mathbf{M}$ is selected as a product of random diagonal matrices (with entries $\pm1$) and a Hadamard matrix. This selection allows multiplication to be performed in $O(d \log d)$ time complexity instead of $O(d^2)$, and can also be merged with the rotary embedding operation without added computational expense.

\end{enumerate}

Our empirical evaluation in \cref{sec:numerical_error} demonstrates that employing these two strategies can reduce numerical errors by up to a factor of 2.6.

\section{Empirical Validation}
\label{sec:exp}

Our implementation of \textsc{FlashAttention-3} leverages the WGMMA and TMA abstractions provided by CUTLASS~\citep{Thakkar_CUTLASS_2023}. Using these primitives, we assess both the performance and the precision of our approach.

\begin{itemize}

\item \textbf{Benchmarking attention.} We evaluate the performance of \textsc{FlashAttention-3} by recording its runtime across varying sequence lengths and perform a comparative analysis with the baseline PyTorch implementation, \textsc{FlashAttention-2}, the Triton-based version of \textsc{FlashAttention-2} that leverages H100-specific operations, as well as a customized vendor edition of \textsc{FlashAttention-2} from cuDNN tuned for H100 GPUs. Our experiments demonstrate that \textsc{FlashAttention-3} delivers speed improvements of up to 2.0$\times$ over \textsc{FlashAttention-2} and up to 1.5$\times$ faster than the Triton adaptation. Furthermore, \textsc{FlashAttention-3} attains a throughput of up to 740 TFLOPs/s, corresponding to 75\% of the peak theoretical performance on H100 GPUs.
\item \textbf{Ablation study.} We verify through controlled experiments that enhancements such as warp-specialization and pipelined GEMM-softmax underlie the observed acceleration in \textsc{FlashAttention-3}.
\item \textbf{Accuracy of FP8 attention.} Our assessments indicate that by employing block quantization and incoherent computation, the FP8 \textsc{FlashAttention-3} achieves a 2.6$\times$ reduction in numerical inaccuracies.

\subsection{Benchmarking Attention}
\label{subsec:benchmark_attn}

We evaluate the performance of various attention mechanisms using an H100 80GB SXM5 GPU under multiple configurations, specifically considering cases with and without a causal mask and utilizing head dimensions of 64 and 128, all with FP16 input representations. The performance outcomes are presented in~\cref{fig:benchmark_attn_fwd} and~\cref{fig:benchmark_attn_bwd}. The data illustrate that \textsc{FlashAttention-3} achieves approximately 1.5-2.0$\times$ higher speed than \textsc{FlashAttention-2} in the forward computation, and 1.5-1.75$\times$ faster in the backward computation. Relative to conventional attention implementations, \textsc{FlashAttention-3} delivers between 3 and 16 times speedup. Notably, for sequence lengths of 1,000 or more, \textsc{FlashAttention-3} even outperforms a proprietary vendor library (cuDNN, which is closed-source) tailored for optimal use on H100 GPUs.

\paragraph{Benchmark settings:} We vary the sequence length as 512, 1k, ..., 16k, and set batch size so that the total number of tokens is 16k. We set the hidden dimension to 2048, and head dimension to be either 64, 128, or 256 (i.e., 32 heads, 16 heads, or 8 heads). To calculate the FLOPs of the forward pass, we use:

\begin{equation*}
  4 \cdot \text{seqlen}^2 \cdot \text{head dimension} \cdot \text{number of heads}.
\end{equation*}

Employing causal masking requires dividing this value by 2, reflecting that roughly 50% of the entries are computed. For the backward pass, the number of FLOPs is obtained by scaling the forward pass FLOPs by a factor of 2.5—this arises because the forward pass consists of 2 matrix multiplications, while the backward pass involves 5 matrix multiplications as a result of recomputation.

Additionally, we evaluate the forward pass runtime using FP8 within comparable conditions. The findings for headdim 256 are presented in ~\cref{fig:benchmark_attn_fp8}, while comprehensive results can be found in \cref{sec:benchmark_attn_fp8_full}.

\subsection{Ablation Study: 2-Stage Pipelining Experiments}
\label{sec:2-stage-pipelining-experiments}

We perform an ablation study on both the 2-stage WGMMA-softmax pipelining and warp-specialization techniques for non-causal FP16 \textsc{FlashAttention-3}, maintaining fixed parameters $\{\text{batch}, \text{seqlen}, \text{nheads}, \text{hdim}\} = \{ 4, 8448, 16, 128\}$. As demonstrated in~\cref{table:ablation_pipelining}, our methodological enhancements—namely, introducing asynchrony via warp-specialization and enabling overlap between GEMM and softmax computations—yield a notable performance increase, raising throughput from 570 to 661 TFLOPs.

\subsection{Numerical Error Validation}
\label{sec:numerical_error}

Given the recent focus on the numerical error associated with \textsc{FlashAttention}~\citep{golden2024flash}, we conduct a comparison involving \textsc{FlashAttention-2}, \textsc{FlashAttention-3}, as well as a conventional attention implementation, benchmarking all against a high-precision FP64 reference. To emulate the presence of outlier features and activations commonly observed in LLMs~\citep{dettmers2208llm, sun2024massive}, we construct $\mathbf{Q}, \mathbf{K}, \mathbf{V}$ using the distribution outlined below:

\begin{equation*}
  \mathcal{N}(0, 1) + \mathcal{N}(0, 100) \cdot \mathrm{Bernoulli}(0.001).
\end{equation*}

In this setup, the majority of entries follow a normal distribution with a mean of zero and a standard deviation of 1, while an independent, normally distributed perturbation with a standard deviation of 10 is added to 0.1\% of the entries. The root mean squared error (RMSE) results are reported in~\cref{table:numerical_error}. When using FP16, both \textsc{FlashAttention-2} and \textsc{FlashAttention-3} yield RMSE values that are 1.7$\times$ smaller than those produced by the standard approach, due to retaining intermediate softmax computations in FP32. For the baseline attention method in FP8, per-tensor scaling is employed, with the matrix multiplication accumulator operating in FP32 and the intermediate softmax results stored in FP16. Owing to the utilization of block quantization and incoherent processing, \textsc{FlashAttention-3} in FP8 achieves a 2.6$\times$ improvement in accuracy relative to this baseline.

\section{Dicussion, Limitations, Conclusion}
\label{sec:discussion}

Through the development of \textsc{FlashAttention-3}, our findings illustrate that leveraging innovative programming methods and advances in hardware—particularly asynchrony and low-precision computation—can substantially enhance both the performance and precision of attention mechanisms. Our approach achieves a 1.5-2.0$\times$ speedup over \textsc{FlashAttention-2} and reduces FP8 numerical errors by a factor of 2.6 in comparison to conventional per-tensor quantization techniques. Nevertheless, several challenges remain for future work: improving support for LLM inference, adopting a persistent kernel architecture for the FP8 kernel,\footnote{Our benchmarks use a persistent kernel and load-balancing strategy in FP16 \textsc{FlashAttention-3}, whereas FP8 \textsc{FlashAttention-3} lacks these features, which partly accounts for its lower performance with short sequences and causal masking relative to the FP8 cuDNN kernels.} and systematically evaluating the impact of low-precision computations on large-scale model training. Although this study concentrated on Hopper GPUs, we anticipate that the underlying approaches can generalize to a range of hardware accelerators. It is our expectation that by refining the efficiency and accuracy of the attention primitive, further progress will be enabled in applications requiring extended contextual understanding.

\end{document}